% Phase 5: PPO optimisation diagnostic
% Auto-authored to accompany phase5_ppo_optimisation_report.md
% Training-env seed 70000, learner seeds 0-2, 64,000 transitions (groups A/B/C),
% checkpoints at t=0,4000,8000,16000,24000,32000,48000,64000.
% Analytical policy and supervised clone used as DIAGNOSTIC REFERENCES only --
% not assumed optimal; primary success criteria are objective, cross-seed
% variance, catastrophic-policy frequency, and economically sensible behaviour.

\subsection*{Section 2: clone-to-actor weight transfer verification}
\begin{table}[h]
\centering
\begin{tabular}{lr}
\hline
Metric & Value \\
\hline
Linear layers copied (of 3: 2 hidden + output) & 3 / 3 (all seeds) \\
Max abs.\ deterministic-output diff, PPO actor vs.\ clone (dense 9,471-state grid) & $4.17\times10^{-7}$ \\
Critic touched by transfer & No (verified: independent per learner seed) \\
log\_std touched by transfer & No (SB3 default log\_std\_init=0.0 preserved) \\
\hline
\end{tabular}
\caption{The clone-to-PPO-actor weight transfer is exact to floating-point
precision; the critic and log\_std are untouched, confirmed by construction
and by two dedicated tests.}
\end{table}

\subsection*{Section 3: initial exploration diagnostic (t=0, clone-initialised, before any PPO update)}
\begin{table}[h]
\centering
\begin{tabular}{lrrrr}
\hline
Learner seed & log\_std & Deterministic $J$ & Stochastic $J$ & Frac.\ sampled actions near bound \\
\hline
0 & 0.0 & 29.41 & $-$15.15 & 0.612 \\
1 & 0.0 & 29.41 & $+$4.91  & 0.611 \\
2 & 0.0 & 29.41 & $-$7.28  & 0.606 \\
\hline
Mean & -- & 29.41 & $-$5.84 & 0.610 \\
\hline
\end{tabular}
\caption{At t=0, action std = $e^{0.0}=1.0$ on a $[-1,1]$ action range: 61\% of
sampled actions land within 0.05 of a bound. The deterministic clone policy
(J=29.4) is far better than the SAME actor's own stochastic behaviour
(mean J=$-$5.8, two of three seeds net losses) -- before a single PPO update.}
\end{table}

\subsection*{Section 4-5: controlled training, deterministic objective at each checkpoint}
\begin{table}[h]
\centering
\small
\begin{tabular}{lrrrrrrrr}
\hline
Group & Seed & t=0 & t=4k & t=8k & t=16k & t=24k & t=32k & t=48k & t=64k \\
\hline
A (random) & 0 & 20.5 & 21.6 & 26.1 & 29.1 & 28.0 & 30.7 & 26.1 & 29.3 \\
A (random) & 1 & 19.9 & 23.3 & 28.0 & 28.9 & 29.4 & 26.5 & 22.9 & 24.7 \\
A (random) & 2 & 20.0 & 20.9 & 28.3 & 31.8 & 25.5 & 27.2 & 24.3 & 28.8 \\
B (clone)  & 0 & 29.4 & 34.3 & 29.8 & 28.3 & 23.0 & 25.9 & 11.2 & 36.0 \\
B (clone)  & 1 & 29.4 & 32.5 & 30.6 & 35.8 & 37.5 & 31.9 & 19.5 & 32.7 \\
B (clone)  & 2 & 29.4 & 37.9 & 32.1 & 36.1 & 39.1 & 43.9 & 32.0 & 43.7 \\
C (frozen) & all & 29.4 & 29.4 & 29.4 & 29.4 & 29.4 & 29.4 & 29.4 & 29.4 \\
\hline
\end{tabular}
\caption{Group C's deterministic objective is bit-identical at every
checkpoint (29.413456 exactly) -- confirms zero evaluation-side drift, so
every change seen in A/B reflects real PPO updates. Group B ends ABOVE both
its own t=0 clone baseline and every Group A seed at t=64k in all 3 seeds,
despite a shared non-monotonic dip around t=24k--48k. Group A's objective
stays consistently below Group B's throughout, and its actor never
approaches the clone's parameter region (see next table).}
\end{table}

\subsection*{Actor drift from the supervised clone (L2 distance in actor-parameter space)}
\begin{table}[h]
\centering
\begin{tabular}{lrrrr}
\hline
Group & t=0 & t=16k & t=32k & t=64k \\
\hline
A (random init), mean of 3 seeds & 13.74 & 13.78 & 13.79 & 13.86 \\
B (clone init), mean of 3 seeds  & 0.00  & 0.29  & 0.32  & 0.53  \\
\hline
\end{tabular}
\caption{Random initialisation makes essentially zero progress toward the
clone's parameter region over 64,000 transitions (distance moves by
$<$1\% of its starting value); clone-initialised runs drift only modestly.}
\end{table}

\subsection*{Section 6: per-update diagnostics -- exploration noise never shrinks}
\begin{table}[h]
\centering
\begin{tabular}{lrrr}
\hline
Group & log\_std\_mean @ t=4k & @ t=32k & @ t=64k \\
\hline
A (random init), mean of 3 seeds & 0.999 & 1.004 & 1.000 \\
B (clone init), mean of 3 seeds  & 1.003 & 1.000 & 0.987 \\
\hline
\end{tabular}
\caption{With ent\_coef=0.0, log\_std\_mean (action std, both dims) stays
within 2\% of its initial value of 1.0 across all 16 updates / 64,000
transitions, in BOTH groups -- PPO's own gradient does not meaningfully
reduce exploration variance within this budget regardless of initialisation.}
\end{table}

\subsection*{Section 7: critic/advantage diagnostic (clone-initialised, t=0, one 4,000-transition rollout, no update)}
\begin{table}[h]
\centering
\begin{tabular}{lr}
\hline
Metric & Value \\
\hline
Independent GAE recompute vs.\ SB3 rollout buffer, max abs.\ diff & $0.0$ (exact match) \\
Mean reward per transition (std) & $-0.049$ (4.14) \\
Mean critic value (std) & $0.085$ (0.134) \\
Mean Monte-Carlo return (std, undiscounted, truncated at buffer end) & $-4.42$ (70.27) \\
Mean / max abs.\ critic error (value $-$ MC return) & $4.50$ / $246.23$ \\
Advantage mean (std, min, max) & $-0.92$ (15.65, $-128.8$, $+76.7$) \\
Frac.\ top-decile-advantage actions near analytical surface ($<$1.0 depth error) & 0.24 \\
Frac.\ bottom-decile-advantage actions near analytical surface & 0.23 \\
\hline
\end{tabular}
\caption{Independent GAE recomputation matches SB3's rollout buffer exactly
(cross-check passes). The randomly-initialised critic's near-zero
predictions are essentially uninformative against realised returns whose
std (70.3) is dominated by the huge sampled-action variance, not by genuine
value-relevant structure -- advantages are correspondingly extreme
(min $-$129, max $+77$) and only weakly associated with analytically
sensible actions (24\%/23\% for top/bottom deciles -- essentially
indistinguishable from each other, i.e.\ the sign of the advantage carries
little information about whether the action was analytically sensible).}
\end{table}

\subsection*{Section 10, conditional diagnostic E: reduced log\_std\_init (clone-initialised, 3 seeds, 32,000 transitions)}
\begin{table}[h]
\centering
\begin{tabular}{lrrrr}
\hline
Configuration & t=0 stoch.\ $J$ & t=16k stoch.\ $J$ & t=32k stoch.\ $J$ & Mean frac.\ near bound \\
\hline
B: log\_std\_init=0.0 (std=1.0), mean of 3 seeds & $-15.15$ & $-10.29$ & $-4.18$ & $\sim$0.5--0.6 \\
E: log\_std\_init=$-$1.5 (std=0.22), mean of 3 seeds & $+26.60$ & $+28.59$ & $+17.13$ & $\sim$0.02--0.03 \\
\hline
\end{tabular}
\caption{Reducing log\_std\_init from 0.0 to $-$1.5 (holding every other
hyperparameter and the clone-initialised actor identical) closes the
deterministic/stochastic performance gap dramatically: the stochastic
objective moves from deeply negative (net loss) to positive and close to
the deterministic objective, with near-bound sampling essentially
eliminated. Deterministic-objective trajectories over the same 32,000
transitions are comparable in magnitude between the two configurations
(E: 29.4$\to$26.7, lower-variance; B: 29.4$\to$33.9, noisier, not yet
showing the t=48k dip since E was only run to 32,000 transitions) --
this diagnostic confirms exploration variance as the direct, fixable cause
of the huge stochastic-performance gap, without yet demonstrating a clear
acceleration of deterministic-objective improvement over this short horizon.}
\end{table}
